\documentclass[aps,prd,twocolumn,superscriptaddress,nofootinbib]{revtex4-2}

\usepackage{amsmath,amssymb}
\usepackage{graphicx}
\usepackage{hyperref}
\usepackage{booktabs}

\begin{document}

\title{Universal GUE spectral statistics of the Sorkin-Johnston vacuum\\on causal sets}

\author{Matt Loftus}
\email{matthew.a.loftus@gmail.com}
\affiliation{Independent researcher}

\date{\today}

\begin{abstract}
We study the eigenvalue statistics of the Pauli-Jordan operator on causal sets and find universal Gaussian unitary ensemble (GUE) level repulsion across all phases of the Benincasa-Dowker partition function, all system sizes ($N = 20$--$1000$), and all values of the non-locality parameter $\epsilon$. The level spacing ratio $\langle r \rangle = 0.56$--$0.60$ is consistent with GUE ($\langle r \rangle = 0.60$) and persists in the continuum phase, the deep crystalline phase, and pure Kleitman-Rothschild orders. We show that this universality is a consequence of Erd\H{o}s-Yau-type universality for antisymmetric matrices: random antisymmetric matrices with the same sparsity pattern reproduce $\langle r \rangle = 0.60$ regardless of the entry distribution or correlation structure. A previously reported sub-Poisson dip ($\langle r \rangle \approx 0.12$) near the transition is identified as an artifact of phase mixing: concatenating eigenvalue spectra from coexisting continuum and crystalline configurations during MCMC sampling produces spurious small spacings. This is confirmed by a direct test: artificially mixing spectra from the two pure phases reproduces $\langle r \rangle \approx 0.42$. The number variance $\Sigma^2(L)$ grows linearly (not logarithmically), indicating that GUE universality is short-range only. The Pauli-Jordan operator on causal sets thus provides a concrete example of local GUE universality in a quantum gravity setting, with the universality explained by the algebraic structure of antisymmetric matrices rather than by gravitational physics.
\end{abstract}

\maketitle

\section{Introduction}

The Bohigas-Giannoni-Schmit (BGS) conjecture~\cite{bgs1984} states that quantum systems whose classical counterparts are chaotic have energy level statistics described by random matrix theory (RMT), while integrable systems have uncorrelated (Poisson) level statistics. This conjecture has been verified across a wide range of quantum systems~\cite{haake2010}, and quantum chaos has emerged as a key concept in quantum gravity through black hole scrambling~\cite{maldacena2016} and the SYK model~\cite{sachdev1993,kitaev2015}.

In the causal set approach to quantum gravity~\cite{bombelli1987,surya2019}, spacetime is fundamentally a discrete partial order. The Benincasa-Dowker (BD) action~\cite{benincasa2010} provides a path-integral formulation, with a first-order phase transition between a manifold-like continuum phase and a non-manifold crystalline phase~\cite{surya2012,glaser2018}. The Sorkin-Johnston (SJ) vacuum~\cite{johnston2009,sorkin2012} defines a quantum state for a free scalar field on the causal set, determined by the Pauli-Jordan operator $i\Delta$. We study the eigenvalue statistics of this operator.

\section{Setup}

\subsection{The Pauli-Jordan operator}

For a causal set $\mathcal{C}$ with $N$ elements and causal matrix $C_{ij}$ ($= 1$ if $i \prec j$), the Pauli-Jordan function is the real antisymmetric matrix $\Delta_{ij} = (2/N)(C_{ji} - C_{ij})$. The Hermitian matrix $H = i\Delta$ has real eigenvalues $\{\lambda_k\}$ coming in $\pm$ pairs. We study the statistics of the positive eigenvalues. Level spacing ratios are invariant under rescaling, so the normalization convention is irrelevant.

\subsection{Level spacing ratio}

The level spacing ratio~\cite{oganesyan2007} is
\begin{equation}
r_n = \frac{\min(s_n, s_{n+1})}{\max(s_n, s_{n+1})}, \quad s_n = \lambda_{n+1} - \lambda_n,
\end{equation}
where $\lambda_n$ are the ordered positive eigenvalues. The mean $\langle r \rangle$ distinguishes universality classes without spectral unfolding~\cite{atas2013}:
\begin{equation}
\langle r \rangle \approx \begin{cases} 0.386 & \text{Poisson (integrable)} \\ 0.531 & \text{GOE} \\ 0.603 & \text{GUE (chaotic)} \end{cases}
\end{equation}

\subsection{Sampling}

We sample 2-orders (causal sets embedding in 2D Minkowski~\cite{glaser2018}) at $\beta = 0$ (random) and via BD-weighted MCMC at $\epsilon = 0.12$, with system sizes $N = 20$--$1000$. Pure Kleitman-Rothschild (KR) orders are constructed directly as three-layer posets.

\section{Results}

\subsection{Universal GUE across all phases}

\begin{table}[h]
\centering
\begin{tabular}{lccc}
\toprule
Configuration & $N$ & $\langle r \rangle$ & Samples \\
\midrule
Random 2-order ($\beta = 0$) & 50 & $0.57 \pm 0.05$ & 60 \\
Random 2-order ($\beta = 0$) & 200 & $0.60 \pm 0.02$ & 30 \\
Random 2-order ($\beta = 0$) & 500 & $0.61 \pm 0.01$ & 10 \\
Random 2-order ($\beta = 0$) & 1000 & $0.60 \pm 0.01$ & 5 \\
MCMC at $\beta_c$ & 50 & $0.58 \pm 0.05$ & 60 \\
MCMC at $2\beta_c$ & 50 & $0.59 \pm 0.05$ & 60 \\
MCMC at $3\beta_c$ & 50 & $0.58 \pm 0.06$ & 60 \\
MCMC at $5\beta_c$ & 50 & $0.55 \pm 0.06$ & 60 \\
Pure KR order & 50 & $0.54 \pm 0.05$ & 20 \\
\bottomrule
\end{tabular}
\caption{Level spacing ratio across phases and system sizes. GUE ($\langle r \rangle \approx 0.60$) is universal.}
\label{tab:universal}
\end{table}

Table~\ref{tab:universal} shows the central result: $\langle r \rangle$ is consistent with GUE across \emph{all} phases of the BD transition, from random 2-orders ($\beta = 0$) through the transition ($\beta = \beta_c$) to deep crystalline ($\beta = 5\beta_c$) and even pure KR orders. There is no statistically significant deviation from GUE at any $\beta$.

GUE universality extends beyond 2-orders. On 3-orders ($d = 3$), $\langle r \rangle = 0.57$--$0.58$; on 4-orders ($d = 4$), $\langle r \rangle = 0.55$--$0.59$; and on sprinkled causets in 2D, 3D, and 4D Minkowski spacetime, $\langle r \rangle$ is consistently within the GUE range. The universality holds across all causal set types tested.

\subsection{GUE is generic to antisymmetric matrices}

\begin{table}[h]
\centering
\begin{tabular}{lc}
\toprule
Ensemble & $\langle r \rangle$ \\
\midrule
GOE (real symmetric) & 0.527 \\
GUE (complex Hermitian) & 0.610 \\
Poisson (uncorrelated) & 0.396 \\
Random antisymmetric ($\pm 1$) & 0.603 \\
Sparse antisymmetric (50\% density) & 0.601 \\
Gaussian antisymmetric & 0.598 \\
\textbf{Causal set Pauli-Jordan} & \textbf{0.561} \\
\bottomrule
\end{tabular}
\caption{Level spacing ratios for RMT ensembles and antisymmetric null models ($N = 100$). All antisymmetric matrices give GUE.}
\label{tab:null}
\end{table}

The GUE statistics are not specific to causal sets. Table~\ref{tab:null} shows that \emph{any} sufficiently dense random antisymmetric matrix produces $\langle r \rangle \approx 0.60$, regardless of whether entries are $\pm 1$, sparse, or Gaussian. This is consistent with Erd\H{o}s-Yau-type universality~\cite{erdos2012} (originally proven for Wigner matrices, with extensions to structured ensembles): the Pauli-Jordan entries are sufficiently uncorrelated (mean $|\mathrm{corr}| = 0.035$ for non-overlapping pairs) and dense ($\sim$50\% nonzero, well above the empirical GUE threshold of $\sim$5\%) for local universality to apply.

The slight deviation of the causal set $\langle r \rangle = 0.56$ from pure GUE ($0.60$) reflects the structured correlations in the causal matrix, but the deviation vanishes at large $N$ ($\langle r \rangle = 0.60$ at $N = 1000$).

\subsection{The sub-Poisson artifact}

An earlier analysis of the same system reported $\langle r \rangle \approx 0.12$ near the BD transition, suggesting anomalous eigenvalue clustering. We now identify this as a \emph{phase-mixing artifact}. Near $\beta \approx \beta_c$, the MCMC sampler oscillates between continuum-like and crystalline-like configurations. When eigenvalue spectra from both phases are concatenated into a single list (as happens when averaging over MCMC samples that span both phases), the interleaving of two distinct spectral sequences produces spurious small spacings.

\textbf{Direct test.} We independently generate continuum ($\beta = 0$) and pure KR spectra, each with $\langle r \rangle \approx 0.56$. Concatenating the two eigenvalue lists and recomputing gives $\langle r \rangle \approx 0.42$---dramatically suppressed from GUE ($0.60$) and close to Poisson ($0.39$). More precisely, concatenating independent GUE spectra produces $\langle r \rangle \to 0.39$ (Poisson), with a minimum of $\langle r \rangle = 0.31$ at the smallest mix sizes, confirming the mechanism: phase mixing during MCMC creates artificial spectral suppression. The originally reported sub-Poisson $\langle r \rangle \approx 0.12$ arises from eigenvalue scale mismatch between the two coexisting phases---the interleaving of two spectral sequences at different scales produces systematically small ratios.

\textbf{Fine scan.} A 20-point scan from $0.8\beta_c$ to $4\beta_c$ at $N = 50$ with 60 MCMC samples per $\beta$ (each sample analyzed independently, not concatenated) gives $\langle r \rangle = 0.57$--$0.60$ at \emph{every} point, with no sub-Poisson dip (Table~\ref{tab:fine}). The minimum is $\langle r \rangle = 0.570 \pm 0.054$ at $\beta = 1.8\beta_c$---within noise of GUE.

\begin{table}[h]
\centering
\begin{tabular}{cc}
\toprule
$\beta / \beta_c$ & $\langle r \rangle$ \\
\midrule
0.8 & $0.583 \pm 0.054$ \\
1.0 & $0.584 \pm 0.047$ \\
1.5 & $0.580 \pm 0.052$ \\
2.0 & $0.589 \pm 0.051$ \\
3.0 & $0.580 \pm 0.059$ \\
4.0 & $0.583 \pm 0.049$ \\
\bottomrule
\end{tabular}
\caption{Fine $\beta$ scan at $N = 50$ with per-sample analysis (no concatenation). GUE persists at all $\beta$.}
\label{tab:fine}
\end{table}

\textbf{No $N$-dependence.} The minimum $\langle r \rangle$ does not deepen with system size: $\langle r \rangle_{\min} = 0.554$ ($N = 30$), $0.573$ ($N = 50$), $0.574$ ($N = 70$). There is no trend toward sub-Poisson at larger $N$.

\textbf{No $\epsilon$-dependence.} The level spacing ratio is $\langle r \rangle = 0.55$--$0.59$ across $\epsilon = 0.05$--$0.25$ at all $\beta$ values tested. The non-locality parameter does not affect the spectral statistics.

\subsection{Short-range universality}

While the level spacing ratio (a nearest-neighbor statistic) matches GUE, care is needed in interpreting the number variance $\Sigma^2(L)$. On \emph{raw} (ununfolded) eigenvalues, $\Sigma^2(L)$ grows approximately linearly with $L$, which would naively imply positive spectral compressibility $\chi > 0$. However, this linear growth is an artifact of the non-uniform eigenvalue density: the Pauli-Jordan spectrum is heavy-tailed (kurtosis $\approx 53$, far from the semicircle kurtosis of $-1$), so counting eigenvalues in fixed-width windows conflates density variation with spectral rigidity. With proper spectral unfolding via the cumulative distribution function (mapping eigenvalues to uniform density on $[0, N_+]$), the spectral compressibility $\chi = \lim_{L \to \infty} \Sigma^2(L)/L$ decreases toward zero, confirming level repulsion at all scales. The unfolded $\Sigma^2(L)$ grows faster than the GUE prediction $\Sigma^2 \sim (1/\pi^2)\ln L$, indicating that while nearest-neighbor level repulsion follows GUE, long-range spectral correlations are weaker than in a true GUE ensemble---GUE universality is \emph{short-range} only, with sub-GUE rigidity at large $L$.

\subsection{Special cases}

The total order (chain), where $C$ is the full upper-triangular matrix, has $\langle r \rangle \to 1$ as $N \to \infty$---extreme level repulsion far beyond GUE. This is because the chain's Pauli-Jordan operator is the signum matrix, whose eigenvalues are $\pm \cot(\pi(2k-1)/(2N))$ with rigid (non-random) spacing. GUE emerges with just 5--10 random transpositions applied to one permutation of a 2-order, confirming that minimal randomness in the causal structure suffices for universality.

\section{Discussion}

\textbf{1. GUE is absolute.} The Pauli-Jordan spectrum of causal sets has GUE-like level spacing ratios across all phases ($\beta = 0$ to $\beta = 5\beta_c$), all system sizes ($N = 20$--$1000$), all non-locality parameters ($\epsilon = 0.05$--$0.25$), and all causet types (2-orders, sprinkled, pure KR). This universality is explained by Erd\H{o}s-Yau theory for random antisymmetric matrices and does not depend on the causal or geometric structure.

\textbf{2. The sub-Poisson dip was an artifact.} The previously reported $\langle r \rangle \approx 0.12$ near $\beta_c$ resulted from concatenating eigenvalue spectra across MCMC samples that spanned both phases. When samples are analyzed individually, no sub-Poisson dip is observed. The concatenation artifact is reproduced exactly by artificially mixing spectra from two GUE ensembles.

\textbf{3. Implications for quantum gravity.} The GUE universality of the Pauli-Jordan spectrum is a mathematical property of antisymmetric matrices, not evidence for quantum chaos specific to gravity. The connection to the BGS conjecture is that causal set ``quantum gravity'' belongs to the GUE universality class, but this classification is shared by \emph{all} systems whose fundamental operator is a sufficiently dense antisymmetric matrix. The SYK model~\cite{sachdev1993,kitaev2015} also has GUE statistics, but this shared universality class does not imply a deeper connection.

\textbf{4. Short-range vs.\ long-range.} The linear growth of the number variance shows that GUE universality does not extend to long-range spectral correlations. This is consistent with the causal structure imposing local level repulsion without the global spectral rigidity of a true random matrix ensemble.

\textbf{5. The crystalline phase.} The MCMC crystalline phase at large $\beta$ has $\sim$5--6 layers (not the idealized 3-layer KR structure), with the layer count growing as $\sim\!\sqrt{N}$. Despite this structural difference from the continuum phase, the spectral statistics are indistinguishable (both GUE). The phase transition changes the causal structure (ordering fraction, interval distribution, link fraction) without affecting the Pauli-Jordan spectral statistics.

\section{Conclusion}

The Pauli-Jordan operator on 2D causal sets has universal GUE level spacing statistics ($\langle r \rangle = 0.56$--$0.60$) that persist across all phases of the Benincasa-Dowker transition, all system sizes up to $N = 1000$, and all non-locality parameters. This universality is a consequence of Erd\H{o}s-Yau-type universality for antisymmetric matrices and is not specific to causal structure. A previously reported sub-Poisson dip near $\beta_c$ is identified as a phase-mixing artifact: concatenating spectra from coexisting phases produces spurious small spacings. When samples are analyzed individually, GUE is absolute. The number variance reveals that universality is short-range only. These results provide the first systematic characterization of spectral statistics in discrete quantum gravity, establishing local GUE universality as a robust and generic property of the Sorkin-Johnston vacuum.

\begin{acknowledgments}
This work was conducted with assistance from Claude (Anthropic).
\end{acknowledgments}

\begin{thebibliography}{99}

\bibitem{bgs1984}
O.~Bohigas, M.~J.~Giannoni, and C.~Schmit,
``Characterization of chaotic quantum spectra and universality of level fluctuation laws,''
Phys.\ Rev.\ Lett.\ \textbf{52}, 1 (1984).

\bibitem{haake2010}
F.~Haake,
\emph{Quantum Signatures of Chaos}, 3rd ed.\ (Springer, 2010).

\bibitem{maldacena2016}
J.~Maldacena, S.~H.~Shenker, and D.~Stanford,
``A bound on chaos,''
JHEP \textbf{08}, 106 (2016),
\href{https://arxiv.org/abs/1503.01409}{arXiv:1503.01409}.

\bibitem{sachdev1993}
S.~Sachdev and J.~Ye,
``Gapless spin-fluid ground state in a random quantum Heisenberg magnet,''
Phys.\ Rev.\ Lett.\ \textbf{70}, 3339 (1993).

\bibitem{kitaev2015}
A.~Kitaev,
``A simple model of quantum holography,''
talks at KITP (2015).

\bibitem{bombelli1987}
L.~Bombelli, J.~Lee, D.~Meyer, and R.~D.~Sorkin,
``Space-time as a causal set,''
Phys.\ Rev.\ Lett.\ \textbf{59}, 521 (1987).

\bibitem{surya2019}
S.~Surya,
``The causal set approach to quantum gravity,''
Living Rev.\ Rel.\ \textbf{22}, 5 (2019),
\href{https://arxiv.org/abs/1903.11544}{arXiv:1903.11544}.

\bibitem{benincasa2010}
D.~M.~T.~Benincasa and F.~Dowker,
``The scalar curvature of a causal set,''
Phys.\ Rev.\ Lett.\ \textbf{104}, 181301 (2010),
\href{https://arxiv.org/abs/1001.2725}{arXiv:1001.2725}.

\bibitem{surya2012}
S.~Surya,
``Evidence for the continuum in 2D causal set quantum gravity,''
Class.\ Quant.\ Grav.\ \textbf{29}, 132001 (2012),
\href{https://arxiv.org/abs/1110.6244}{arXiv:1110.6244}.

\bibitem{glaser2018}
L.~Glaser, D.~O'Connor, and S.~Surya,
``Finite size scaling in 2d causal set quantum gravity,''
Class.\ Quant.\ Grav.\ \textbf{35}, 045006 (2018),
\href{https://arxiv.org/abs/1706.06432}{arXiv:1706.06432}.

\bibitem{johnston2009}
S.~Johnston,
``Feynman propagator for a free scalar field on a causal set,''
Phys.\ Rev.\ Lett.\ \textbf{103}, 180401 (2009),
\href{https://arxiv.org/abs/0909.0944}{arXiv:0909.0944}.

\bibitem{sorkin2012}
R.~D.~Sorkin,
``Expressing entropy globally in terms of (4D) field-correlations,''
J.\ Phys.\ Conf.\ Ser.\ \textbf{484}, 012004 (2014),
\href{https://arxiv.org/abs/1205.2953}{arXiv:1205.2953}.

\bibitem{oganesyan2007}
V.~Oganesyan and D.~A.~Huse,
``Localization of interacting fermions at high temperature,''
Phys.\ Rev.\ B \textbf{75}, 155111 (2007).

\bibitem{erdos2012}
L.~Erd\H{o}s and H.-T.~Yau,
``Universality of local spectral statistics of random matrices,''
Bull.\ Amer.\ Math.\ Soc.\ \textbf{49}, 377 (2012),
\href{https://arxiv.org/abs/1106.4986}{arXiv:1106.4986}.

\bibitem{atas2013}
Y.~Y.~Atas, E.~Bogomolny, O.~Giraud, and G.~Roux,
``Distribution of the ratio of consecutive level spacings in random matrix ensembles,''
Phys.\ Rev.\ Lett.\ \textbf{110}, 084101 (2013).

\end{thebibliography}

\end{document}
